\section{Boundary divisors and canonical factor removal}\label{sec:offcentral-obstructions}

This appendix studies the boundary function \(r_\alpha(t)=-\ord_{\alpha+it}M\) when the comparison abscissa lies strictly to the right of \(1/2\). It proves the two assertions in Theorem~\ref{thm:boundary-method} and gives explicit examples delimiting their scope. Throughout, positive definiteness is on \(\R_d\), while local finiteness refers to the ordinary topology on \(\R\). A generalized Dirichlet series has distinct positive frequencies that are locally finite in \([0,\infty)\), and converges absolutely in some right half-plane.

The canonical factor of a boundary divisor means the finite product of exponential and linear factors in \eqref{eq:ob3-9}, normalized so that its logarithmic derivative has no constant term at \(+\infty\). The results concern this precise operation. They do not replace the nonvanishing condition \eqref{eq:RB}, and do not narrow the remaining interval for \(\alpha\).

\subsection{Simple boundary profiles}

\begin{proposition}[Simple boundary profiles]\label{prop:ob-simple}

Let \(r:\mathbb R_d\to\mathbb Z\) be positive definite, with \(r(0)=1\), and suppose its support is locally finite in the ordinary real topology. Exactly one of the following forms occurs:

\begin{equation}\label{eq:ob2-1}
r=\mathbf1_{\{0\}};
\end{equation}

\begin{equation}\label{eq:ob2-2}
r=\mathbf1_{h\mathbb Z}\quad\text{for some }h>0;
\end{equation}

\begin{equation}\label{eq:ob2-3}
r(nh)=(-1)^n\quad(n\in\mathbb Z),\qquad
r(t)=0\quad(t\notin h\mathbb Z),
\end{equation}

again for some \(h>0\). No idempotent theorem is needed for this special case.

\end{proposition}

\begin{proof} Positivity gives \(|r(t)|\le1\), so the values are \(-1,0,1\). Real-valuedness also gives \(r(-t)=r(t)\). Realize the kernel \(r(t-u)\) as a Gram kernel of unit vectors \(v_t\). This realization can be obtained directly by completing the space of finite formal linear combinations modulo its null space; it invokes no continuity of \(r\).

If \(|r(t)|=1\), equality in Cauchy--Schwarz gives \(v_t=r(t)v_0\). Translation is an isometry of the Gram construction, and therefore

\begin{equation}\label{eq:ob2-4}
r(t+u)=r(t)r(u)\qquad(u\in\mathbb R).
\end{equation}

Consequently \(H=\{t:|r(t)|=1\}\) is a subgroup and \(r|_H\) is a character with values in \(\{1,-1\}\). Since the other values of \(r\) are zero, \(H\) is its support. A locally finite subgroup of \(\mathbb R\) is either \(\{0\}\) or \(h\mathbb Z\): in the nonzero case its positive elements have a least member \(h\), and Euclidean division reduces every member modulo \(h\). A character of \(h\mathbb Z\) into \(\{1,-1\}\) is determined by its value at \(h\). These are precisely \eqref{eq:ob2-1}--\eqref{eq:ob2-3}. Conversely all three are positive definite, by grouping finite Gram sums according to cosets of \(H\). \end{proof}

\paragraph{Arithmetic interpretation.}

If the entire quadratic Hecke function \(B=L(s,\eta_E)L(s,\rho_5)\) has a \emph{simple} zero at a real comparison pole \(\alpha\), then \(m=1\). Indeed, \(\Lambda(s,\pi)\) is entire at \(\alpha\), so a pole of \(M=\Lambda(\pi)/\Lambda(\rho_5)\) forces a positive zero order of \(L(s,\rho_5)\). The entire factor \(L(s,\eta_E)\) has nonnegative zero order there. Their sum is one, so \(L(s,\rho_5)\) has order one, \(L(s,\eta_E)\) has order zero, and the numerator has order zero. This implication is conditional on simplicity; simplicity of an arbitrary off-central zero of \(B\) is not asserted.

\subsection{Canonical boundary removal}

Here ``canonical removal'' means dividing by the finite-order elementary factor having exactly the boundary divisor in Proposition~\ref{prop:ob-simple}, with no extra divisors. The harmless freedom to multiply that factor by \(e^{as+b}\) only changes a conductor constant in its logarithmic derivative. We use the Euler normalization in which that constant is zero. Arbitrary zero-free analytic factors are not permitted: allowing them would simply move the unknown part of \(Q_\Omega\) into the factor being removed.

\begin{proposition}[Removal of a simple boundary divisor]\label{prop:ob-simple-peeling}

Suppose \(\alpha>0\), and \(D_\Omega\) is as in \eqref{eq:Dpositive}, absolutely convergent on a right half-plane. For each of the three simple-pole profiles in Proposition~\ref{prop:ob-simple}, let \(P_\alpha\) be the canonical factor below. Then

\begin{equation}\label{eq:ob3-1}
-\frac{(Q_\Omega/P_\alpha)'}{Q_\Omega/P_\alpha}
=D_\Omega-D_{P_\alpha}
\end{equation}

cannot be an absolutely convergent generalized Dirichlet series with nonnegative coefficients and positive frequencies. The same conclusion holds if the bounded finite bad-prime contributions in \eqref{eq:Dar} are included.

\end{proposition}

\begin{proof}\textit{Singleton profile.} Take

\begin{equation}\label{eq:ob3-2}
P_\alpha(s)=(s-\alpha)^{-1},\qquad
D_{P_\alpha}(s)=\frac1{s-\alpha}.
\end{equation}

Any absolutely convergent prime-power series in \eqref{eq:Dpositive} or \eqref{eq:Dar} is \(O(e^{-c\sigma})\) as real \(\sigma\to+\infty\), for some \(c>0\). Therefore

\begin{equation}\label{eq:ob3-3}
D_\Omega(\sigma)-\frac1{\sigma-\alpha}<0
\end{equation}

for all sufficiently large \(\sigma\). This is incompatible with a nonnegative Dirichlet series. In Laplace-transform language, \(1/(s-\alpha)\) corresponds to the continuous positive density \(e^{\alpha x}\,dx\) on \(x>0\), whereas \(D_\Omega\) comes from a positive atomic measure. Its removal subtracts that continuous density. A conductor constant at frequency zero does not remove the negative continuous part.

\textit{Lattice profile.} If \(r=\mathbf1_{h\mathbb Z}\), put \(\lambda=2\pi/h\) and

\begin{equation}\label{eq:ob3-4}
P_\alpha(s)=(1-e^{-\lambda(s-\alpha)})^{-1}.
\end{equation}

Then, for \(\Re s>\alpha\),

\begin{equation}\label{eq:ob3-5}
D_{P_\alpha}(s)=\lambda\sum_{n\ge1}
e^{\alpha n\lambda}e^{-n\lambda s}.
\end{equation}

At frequencies \(n\lambda\), the prime-power series can involve at most one underlying prime. In fact, \(n\lambda=k\log p\) and \(n'\lambda=k'\log q\) imply \(p^{kn'}=q^{k'n}\), hence \(p=q\). The original coefficient on this entire frequency progression is consequently bounded by \(4\log p\), or by \(10\log p\) when finite bad-prime terms are included; it is zero when no prime occurs. The subtracted coefficient \(\lambda e^{\alpha n\lambda}\) tends to infinity. Thus the remainder has negative coefficients at arbitrarily large \(n\lambda\).

\textit{Alternating lattice profile.} If \(r(nh)=(-1)^n\), put \(\lambda=\pi/h\) and

\begin{equation}\label{eq:ob3-6}
P_\alpha(s)=
\frac{1+e^{-\lambda(s-\alpha)}}{1-e^{-\lambda(s-\alpha)}}.
\end{equation}

Its poles are at the even multiples of \(ih\), and its zeros at the odd multiples, giving exactly the required signed divisor. Moreover

\begin{equation}\label{eq:ob3-7}
D_{P_\alpha}(s)=2\lambda
\sum_{\substack{n\ge1\\n\text{ odd}}}
e^{\alpha n\lambda}e^{-n\lambda s}.
\end{equation}

The same bounded-coefficient argument gives negative remainder coefficients at all sufficiently large odd model frequencies that occur in the union of frequency sets.

In the last two cases the union of the lattice frequencies and prime-power logarithms is locally finite and bounded away from zero. Both series converge absolutely sufficiently far to the right. Uniqueness of generalized Dirichlet series applies: if a nonzero coefficient at the least remaining frequency existed in a series representing zero, multiplying by its leading exponential and letting \(\sigma\to+\infty\) would contradict that coefficient. Iterating proves uniqueness. Thus the negative coefficients cannot be replaced by another nonnegative expansion. \end{proof}

The argument also survives canonical removal of the reflected boundary at \(1-\alpha>0\). The reflected divisor is the same in the three cases because the profiles are even. The corresponding positive-frequency logarithmic derivatives are again nonnegative, with \(\alpha\) replaced by \(1-\alpha\), so their subtraction cannot correct the negative terms just exhibited.

\paragraph{Scope of the obstruction.} Proposition~\ref{prop:ob-simple-peeling} excludes an iterative Landau/Cohen argument that requires the canonically peeled remainder to retain a nonnegative Euler logarithmic derivative. It does not exclude a different invariant, a signed explicit formula, or an argument that uses the complete automorphic and Artin structure simultaneously.

\begin{proposition}[Canonical removal at arbitrary multiplicity]\label{prop:ob-peeling}

The conclusion of Proposition~\ref{prop:ob-simple-peeling} holds for every nonzero integer-valued positive-definite boundary profile with locally finite ordinary support, without assuming \(r(0)=1\). More precisely, start with any Cohen decomposition

\begin{equation}\label{eq:ob3-8}
r(t)=\sum_{j=1}^J c_j\mathbf1_{\beta_j+h_j\mathbb Z}(t)
+\sum_{u=1}^U b_u\mathbf1_{\{t_u\}}(t),
\qquad c_j,b_u\in\mathbb Z,\quad h_j>0.
\end{equation}

At a line of positive real part \(\alpha\), its canonical factor is

\begin{equation}\label{eq:ob3-9}
P_\alpha(s)=
\prod_j(1-e^{-\lambda_j(s-\alpha-i\beta_j)})^{-c_j}
\prod_u(s-\alpha-it_u)^{-b_u},
\qquad \lambda_j=2\pi/h_j.
\end{equation}

Then \(D_{P_\alpha}=-P_\alpha'/P_\alpha\) is the Laplace transform of a nonzero nonnegative measure on positive frequencies. Nevertheless \(D_{\rm ar}-D_{P_\alpha}\) cannot have a nonnegative generalized Dirichlet expansion when \(D_{\rm ar}\) satisfies \eqref{eq:Dar} and \eqref{eq:bound}. This is a statement about positivity of the \emph{logarithmic derivative}, not a claim that every possible ordinary coefficient of the quotient is negative.

\end{proposition}

\begin{proof}\textit{Positivity of the canonical factor.} Let \(K=\widehat{\mathbb R_d}\) be the compact group of all characters of the additive discrete real group. Adopt the convention

\begin{equation}\label{eq:ob3-10}
r(t)=\int_K\chi(t)\,d\nu(\chi).
\end{equation}

Bochner's theorem for locally compact abelian groups \citep[Theorem~12.3]{korner2008} supplies a positive finite measure \(\nu\). Let \(m_K\) be Haar probability measure on \(K\). For \(H_j=h_j\mathbb Z\) put

\begin{equation}\label{eq:ob3-11}
K_j=H_j^\perp=\{\chi\in K:\chi(h_j)=1\},
\end{equation}

and let \(m_j\) be Haar probability measure on \(K_j\). Orthogonality gives

\begin{equation}\label{eq:ob3-12}
\widehat{\overline{\chi(\beta_j)}\,m_j}(t)
=\mathbf1_{\beta_j+H_j}(t),\qquad
\widehat{\overline{\chi(t_u)}\,m_K}(t)
=\mathbf1_{\{t_u\}}(t).
\end{equation}

Fourier-transform uniqueness therefore identifies the positive measure in \eqref{eq:ob3-10} as

\begin{equation}\label{eq:ob3-13}
\nu=F_0m_K+\sum_{j=1}^J
c_j\overline{\chi(\beta_j)}\,m_j,
\qquad
F_0(\chi)=\sum_u b_u\overline{\chi(t_u)}.
\end{equation}

Partition the \(h_j\) by rational commensurability. For a class \(I\), choose \(h_I>0\) such that

\begin{equation}\label{eq:ob3-14}
d_j:=h_I/h_j\in\mathbb Z_{>0}\quad(j\in I).
\end{equation}

Then \(H_I=h_I\mathbb Z\) is contained in every \(H_j\) in the class with finite index \([H_j:H_I]=d_j\). Its annihilator \(K_I\) contains each \(K_j\) as a clopen subgroup of index \(d_j\), and

\begin{equation}\label{eq:ob3-15}
m_j=d_j\mathbf1_{K_j}\,m_I.
\end{equation}

The part of \(\nu\) belonging to this class is thus

\begin{equation}\label{eq:ob3-16}
\nu_I=F_I m_I,\qquad
F_I(\chi)=\sum_{j\in I}
c_jd_j\,\overline{\chi(\beta_j)}\,\mathbf1_{K_j}(\chi),
\end{equation}

where \(F_I\) is continuous on \(K_I\).

The measures \(m_K\) and the finitely many \(m_I\) are mutually singular. Indeed, \(m_K(K_I)=0\): the evaluation map \(\chi\mapsto\chi(h_I)\) maps Haar measure to Haar measure on the circle, since \(h_I\) has infinite order in \(\mathbb R_d\). For distinct classes \(I,I'\), the class of \(h_{I'}\) has infinite order in \(\mathbb R_d/H_I\). The same orthogonality argument gives

\begin{equation}\label{eq:ob3-17}
m_I(K_I\cap K_{I'})=0
\quad\text{and}\quad
m_{I'}(K_I\cap K_{I'})=0.
\end{equation}

For example, the pushforward under \(\chi\mapsto\chi(h_{I'})\) has all nonzero Fourier coefficients zero because \(nh_{I'}\notin H_I\) for \(n\ne0\), and hence is circle Haar measure. This proves \eqref{eq:ob3-17} without assuming any metrizability of \(K\).

Since \(\nu\) is positive and these components are mutually singular, \(F_0m_K\) and every \(F_Im_I\) are positive. Haar measure has full support on each compact group, so continuity upgrades almost-everywhere nonnegativity to

\begin{equation}\label{eq:ob3-18}
F_0(\chi)\ge0\quad(\chi\in K),\qquad
F_I(\chi)\ge0\quad(\chi\in K_I).
\end{equation}

Evaluate first at the ordinary characters \(\chi_x(t)=e^{-ixt}\). This gives the nonnegative trigonometric polynomial

\begin{equation}\label{eq:ob3-19}
f_0(x):=F_0(\chi_x)=\sum_u b_u e^{it_ux}\ge0.
\end{equation}

For class \(I\), put \(\lambda_I=2\pi/h_I\). At \(x=n\lambda_I\), the character \(\chi_x\) lies in \(K_I\), and it belongs to \(K_j\) precisely when \(d_j\mid n\). The coefficient of \(e^{-n\lambda_I s}\) in the progression part of \(D_{P_\alpha}\) is therefore

\begin{equation}\label{eq:ob3-20}
e^{\alpha n\lambda_I}u_I(n),\qquad
u_I(n)=\lambda_I F_I(\chi_{n\lambda_I})\ge0.
\end{equation}

The function \(u_I(n)\) is a finite sum of unit-modulus exponentials, because each divisibility indicator is a finite root-of-unity average. Distinct classes share no positive frequency. The partial fractions from the individual points have the integral representation

\begin{equation}\label{eq:ob3-21}
\sum_u\frac{b_u}{s-\alpha-it_u}
=\int_0^\infty e^{-sx}e^{\alpha x}f_0(x)\,dx.
\end{equation}

Equations \eqref{eq:ob3-19}--\eqref{eq:ob3-21} show that

\begin{equation}\label{eq:ob3-22}
D_{P_\alpha}(s)=
\int_0^\infty e^{-sx}e^{\alpha x}f_0(x)\,dx
+\sum_I\sum_{n\ge1}
e^{\alpha n\lambda_I}u_I(n)e^{-n\lambda_I s}
\end{equation}

is a nonnegative Laplace transform for \(\Re s>\alpha\).

\textit{Nontriviality.} If \(f_0\) vanishes identically, linear independence of the distinct real exponentials gives \(F_0=0\). If \(u_I(n)=0\) for every positive integer \(n\), the finite-exponential-sum mean-square argument gives that its coefficients vanish, so it also vanishes for all integers \(n\). The ordinary characters \(\{\chi_{n\lambda_I}:n\in\mathbb Z\}\) are dense in \(K_I\): their annihilator is exactly \(h_I\mathbb Z\), and the double-annihilator theorem \citep[Lemma~12.5.2]{dikranjan2018} identifies their closure. Continuity then gives \(F_I=0\). Thus if every term on the right of \eqref{eq:ob3-22} vanished, \eqref{eq:ob3-13} would give \(\nu=0\), hence \(r=0\). The profile is assumed nonzero, so at least one part is nonzero.

\textit{Failure of positivity after removal.} If \(f_0\ne0\), its nonnegative real-analytic function on the real axis has a first nonzero Taylor term at zero, of even degree \(2q\), with a positive coefficient \(c\). Taylor expansion on a small interval, followed by an exponentially bounded tail, gives

\begin{equation}\label{eq:ob3-23}
\int_0^\infty e^{-(\sigma-\alpha)x}f_0(x)\,dx
\sim \frac{c(2q)!}{(\sigma-\alpha)^{2q+1}}
\quad(\sigma\to+\infty).
\end{equation}

The other terms in \eqref{eq:ob3-22} are nonnegative. In contrast, \(D_{\rm ar}(\sigma)=O(e^{-c_1\sigma})\) for some \(c_1>0\). Therefore the real part of \(D_{\rm ar}(\sigma)-D_{P_\alpha}(\sigma)\) is eventually negative. It cannot be a nonnegative Dirichlet series.

If \(f_0=0\), choose a nonzero class \(I\). A nonzero finite exponential sum \(u_I(n)\ge0\) is at least some \(\epsilon>0\) along an infinite sequence of positive integers: otherwise it would tend to zero, contradicting its positive Cesàro mean square. At these frequencies the subtracted coefficients

\begin{equation}\label{eq:ob3-24}
e^{\alpha n\lambda_I}u_I(n)
\end{equation}

tend to infinity, whereas the coefficient of \(D_{\rm ar}\) on this progression is uniformly bounded by the one-prime argument in Proposition~\ref{prop:ob-simple-peeling}. Hence infinitely many remainder coefficients have negative real part. Uniqueness of generalized Dirichlet series finishes the proof. \end{proof}

The argument also applies after removing both reflected boundaries, because each has a positive canonical logarithmic derivative. It supplies no replacement positivity for the interior divisor. It is a rigorous obstruction to this particular induction scheme, not a proof of the impossibility of other off-central reductions.

\subsection{Rigidity from the complete divisor}

The missing input is not the numerical value \(1/2\) in the coefficient comparison. It is control of the \emph{entire} divisor.

\begin{proposition}[Rigidity from finitely many complete progressions]\label{prop:ob-global}

Assume the comparison representation \eqref{eq:M}, the gamma and reflection equations \eqref{eq:G} and \eqref{eq:FE}, and the coefficient bound \eqref{eq:bound}. Suppose the entire signed divisor \(-\operatorname{ord}M\) is a finite integral combination of vertical arithmetic progressions and finitely many individual points, all lying in \(0<\Re s<1\). Then \(M\) is constant and \(D_{\rm ar}=0\). In particular \(\Omega\) is empty under this exact Euler representation.

\end{proposition}

\begin{proof} Write the progression terms as

\begin{equation}\label{eq:ob4-1}
\sum_{j=1}^J c_j\mathbf1_{
\{\sigma_j+i(\beta_j+h_jn):n\in\mathbb Z\}},
\qquad 0<\sigma_j<1,
\end{equation}

and the point terms as \(\sum_u b_u\mathbf1_{\{\rho_u\}}\), where \(0<\Re\rho_u<1\). Set \(\lambda_j=2\pi/h_j\) and

\begin{equation}\label{eq:ob4-2}
P(s)=\prod_j
(1-e^{-\lambda_j(s-\sigma_j-i\beta_j)})^{-c_j}
\prod_u(s-\rho_u)^{-b_u}.
\end{equation}

This is a meromorphic function of order at most one with the entire divisor of \(M\). Hadamard factorization therefore gives \(M=e^{as+b}P\). In particular

\begin{equation}\label{eq:ob4-3}
-M'/M=-a+R+C,
\quad R(s)=\sum_u\frac{b_u}{s-\rho_u},
\end{equation}

\begin{equation}\label{eq:ob4-4}
C(s)=\sum_j c_j\lambda_j\sum_{n\ge1}
e^{\sigma_j n\lambda_j}e^{i\beta_j n\lambda_j}
e^{-n\lambda_js}.
\end{equation}

Let \(h=-G'/G\). The Euler representation gives

\begin{equation}\label{eq:ob4-5}
-M'/M=-\tfrac12\log A+h+D_{\rm ar}.
\end{equation}

Along the positive real axis, \(R,h\to0\), while \(C,D_{\rm ar}\) are exponentially small. Hence \(a=\tfrac12\log A\) and

\begin{equation}\label{eq:ob4-6}
R(s)-h(s)=D_{\rm ar}(s)-C(s)=O(e^{-cs}).
\end{equation}

The gamma recurrence gives

\begin{equation}\label{eq:ob4-7}
h(s+2)-h(s)=-\frac d s+\frac d{s+1}.
\end{equation}

Consequently the rational function

\begin{equation}\label{eq:ob4-8}
R(s+2)-R(s)+\frac d s-\frac d{s+1}
\end{equation}

is exponentially small on the positive real axis and is therefore identically zero. Its possible poles from \(R(s)\) lie in \(0<\Re s<1\), those from \(R(s+2)\) lie in \(-2<\Re s<-1\), and the two gamma poles lie on the separate lines \(\Re s=0,-1\). Thus \(d=0\), and \(R(s+2)=R(s)\). A rational periodic function that tends to zero at infinity is zero. This leaves

\begin{equation}\label{eq:ob4-9}
D_{\rm ar}=C.
\end{equation}

Partition the finitely many \(\lambda_j\) according to positive rational commensurability. In each class choose a common \(\lambda>0\) such that \(\lambda_j=m_j\lambda\) with positive integers \(m_j\). Different classes never share a nonzero frequency. Within one class the coefficient at \(n\lambda\) has the form

\begin{equation}\label{eq:ob4-10}
\sum_{\sigma\in\Sigma}
e^{\sigma n\lambda}u_\sigma(n),
\end{equation}

where \(\Sigma\subset(0,1)\) is finite and each \(u_\sigma\) is a finite sum of unit-modulus exponentials. Indeed, the restriction \(m_j\mid n\) is a finite average of roots of unity to the \(n\)-th power.

By uniqueness of the generalized Dirichlet expansions, \eqref{eq:ob4-10} equals the coefficient in \(D_{\rm ar}\). That coefficient is uniformly bounded along the progression, because the progression meets powers of at most one prime, as in Proposition~\ref{prop:ob-simple-peeling}. Divide by \(e^{\sigma_{\max}n\lambda}\). All lower exponents tend to zero, and the bounded left side does too; hence \(u_{\sigma_{\max}}(n)\to0\). A finite exponential sum

\begin{equation}\label{eq:ob4-11}
u(n)=\sum_{\ell=1}^Lz_\ell e^{in\theta_\ell}
\end{equation}

with distinct phases modulo \(2\pi\) cannot tend to zero unless every coefficient vanishes: its Cesàro mean square tends to \(\sum|z_\ell|^2\). Remove this largest exponent and repeat through the finite set \(\Sigma\). Every coefficient in every class is zero. Thus \(C=D_{\rm ar}=0\).

Now \eqref{eq:ob4-5} gives \(M(s)=e^{as+b}\). The reflection equation forces \(a=0\), so \(M\) is constant. At a prime in \(\Omega\), \eqref{eq:Dar} has coefficient \(a(p,1)=4\), which is incompatible with \(D_{\rm ar}=0\). \end{proof}

This proposition gives a precise strengthened target for any off-central argument: proving a finite progression description of the \emph{whole} divisor would suffice, even if the progressions lie on several different vertical lines. Boundary positive definiteness alone does not supply that target.

\subsection{Exact prime Euler factors with a boundary lattice}

The following construction preserves much more than an abstract positive Laplace transform. It uses rational primes, exactly the local factors \eqref{eq:Q}, bounded prime-power traces, and ordinary Dirichlet coefficients satisfying a Ramanujan bound. It has a meromorphic continuation through its outer line, whose entire boundary divisor is a single positive lattice comb. It shows concretely why those facts alone do not identify the Euler product with its boundary factor.

\begin{proposition}[An exact prime boundary example]\label{prop:ob-primes}

For every positive integer \(m\), there is an infinite set \(\Omega\) of odd primes such that \(Q_\Omega\) has all of the following properties:

\begin{enumerate}
\def\labelenumi{\arabic{enumi}.}

\item
  Its abscissa of absolute convergence is \(\alpha=3/5\), and \(\sum_{p\in\Omega}p^{-1}<\infty\).
\item
  Its local factors are exactly \(((1+p^{-s})/(1-p^{-s}))^2\), so its normalized logarithmic prime-power coefficients are \(4\) for odd powers and zero for even powers.
\item
  Its ordinary coefficients are nonnegative integers and satisfy \(a(n)\le d_4(n)\ll_\epsilon n^\epsilon\).
\item
  It continues meromorphically to \(\Re s>4/15\), where

\begin{equation}\label{eq:ob5-1}
Q_\Omega(s)=(1-2^{3/5-s})^{-m}e^{G(s)}
\end{equation}

for a holomorphic function \(G\).
\item In that half-plane it has no zeros and has precisely the poles

\begin{equation}\label{eq:ob5-2}
s=\frac35+\frac{2\pi i k}{\log2},\qquad k\in\mathbb Z,
\end{equation}

each of order \(m\).
\item The zero-free factor after removing the displayed boundary factor has logarithmic Dirichlet coefficients

\begin{equation}\label{eq:ob5-3}
-m(\log2)2^{3n/5}
\end{equation}

at every frequency \(n\log2\), and hence does not retain positivity.
\end{enumerate}

\end{proposition}

\begin{proof}\textit{Published input.} \citet[p.~562, final displayed estimate]{bhp2001} prove

\begin{equation}\label{eq:ob5-4}
\pi(x+x^{21/40})-\pi(x)
\ge\frac9{100}\frac{x^{21/40}}{\log x}
\end{equation}

for all sufficiently large \(x\). Only the weak consequence

\begin{equation}\label{eq:ob5-5}
\pi(X+X^{2/3})-\pi(X)\gg X^{2/3}/\log X
\end{equation}

is needed. To see it directly, pack disjoint intervals starting at \(X+j(2X)^{21/40}\); each contains its shorter interval of length \(x^{21/40}\), and there are \(\gg X^{2/3-21/40}\) of them within \([X,X+X^{2/3}]\). Multiplying \eqref{eq:ob5-4} by this number proves \eqref{eq:ob5-5}. The constant \(9/100\) is not needed elsewhere.

\textit{Construction.} Set

\begin{equation}\label{eq:ob5-6}
X_n=2^n,\quad \lambda=\log2,\quad
\alpha=\frac35,\quad\beta=\frac23,
\quad T_n=\frac{m\lambda}{4}X_n^\alpha.
\end{equation}

For all sufficiently large \(n\), the interval

\begin{equation}\label{eq:ob5-7}
I_n=[X_n,X_n+X_n^\beta]
\end{equation}

contains primes of total logarithmic weight \(\gg X_n^\beta>T_n\). The intervals are disjoint. In increasing order, select primes from \(I_n\) until their logarithmic weights first reach \(T_n\). Let \(\Omega_n\) be this finite selection and \(\Omega\) their union, omitting all finitely many initial intervals. All chosen primes are odd. The overshoot is at most the logarithm of one prime, so

\begin{equation}\label{eq:ob5-8}
W_n:=\sum_{p\in\Omega_n}\log p
=T_n+O(\log X_n).
\end{equation}

In particular \(W_n\asymp X_n^\alpha\), and \(\#\Omega_n\asymp X_n^\alpha/\log X_n\). These estimates prove convergence of the prime sum for \(\Re s>\alpha\), divergence at \(\alpha\), and convergence at \(s=1\). The positive Euler product therefore has abscissa exactly \(\alpha\), as does its logarithmic derivative. At \(\alpha\), even \(\sum_{p\in\Omega}p^{-\alpha}\) diverges by comparison with \(\sum_n1/n\).

For any compact set of \(s\)'s, the mean value estimate along the real segment from \(X_n\) to \(p\in I_n\) gives

\begin{equation}\label{eq:ob5-9}
|p^{-s}-X_n^{-s}|
\ll_K X_n^{-\Re s-1}(p-X_n).
\end{equation}

Consequently

\begin{equation}\label{eq:ob5-10}
\sum_{p\in\Omega_n}(\log p)p^{-s}
=T_nX_n^{-s}
+O_K\bigl(X_n^{\alpha+\beta-1-\Re s}
+(\log X_n)X_n^{-\Re s}\bigr).
\end{equation}

The sum of errors converges normally when

\begin{equation}\label{eq:ob5-11}
\Re s>\alpha+\beta-1=\frac4{15}.
\end{equation}

The odd prime powers of exponent at least three converge normally for \(\Re s>\alpha/3=1/5\), since they are bounded on compact sets by a constant times \(\sum_{p\in\Omega}(\log p)p^{-3\Re s}\). Including the finitely many omitted geometric terms, \eqref{eq:ob5-10} therefore yields

\begin{equation}\label{eq:ob5-12}
D_\Omega(s)
=\frac{m\lambda}{e^{\lambda(s-\alpha)}-1}+E(s),
\end{equation}

with \(E\) holomorphic in \(\Re s>4/15\). Initially this identity is an identity of absolutely convergent series in \(\Re s>\alpha\), and its right side supplies the continuation.

Choose a holomorphic primitive \(G\) of \(-E\) on that simply connected half-plane. Since

\begin{equation}\label{eq:ob5-13}
-\frac d{ds}\log(1-e^{-\lambda(s-\alpha)})^{-m}
=\frac{m\lambda}{e^{\lambda(s-\alpha)}-1},
\end{equation}

the quotient of \(Q_\Omega\) and the right side of \eqref{eq:ob5-1} is constant in the original half-plane. Choose the additive constant of \(G\) to make this constant one. This proves \eqref{eq:ob5-1}--\eqref{eq:ob5-2}, including all divisor orders.

The local expansion is

\begin{equation}\label{eq:ob5-14}
\left(\frac{1+x}{1-x}\right)^2
=1+4\sum_{k\ge1}k x^k.
\end{equation}

Since \(4k\le\binom{k+3}{3}=d_4(p^k)\), multiplicativity gives the claimed ordinary coefficient bound. Finally, \(2\notin\Omega\), so the original logarithmic derivative has no term at any \(n\log2\). Subtracting the model derivative in \eqref{eq:ob5-12} gives exactly \eqref{eq:ob5-3}. \end{proof}

\paragraph{Scope of the example.}

There is no claim that \(\Omega\) comes from an \(S_5\) representation, and no claim of global meromorphic continuation or order at most one. In fact the reflection functional equation required in \eqref{eq:FE} \emph{fails already on the displayed continuation domain}: the real pole at \(3/5\) would require a pole at \(2/5\), whereas \eqref{eq:ob5-1} is holomorphic and nonzero at \(2/5>4/15\). Multiplication by a nonvanishing conductor exponential cannot repair that missing reflected divisor.

Thus Proposition~\ref{prop:ob-primes} is a counterexample to a boundary-only extension, not to the complete arithmetic conjecture. It establishes that exact prime Euler factors, the Ramanujan coefficient bound, integral boundary orders, Cohen's lattice conclusion, and continuation a fixed distance across the boundary do not by themselves force a finite exceptional set.

\subsection{A symmetric order-one example without an interior gap}

Local finiteness of a divisor gives a horizontal gap on every bounded-height compact set after one removes its boundary points. It does not give a gap uniform in height, even for a globally meromorphic function of order one with the exact reflection symmetry.

\begin{proposition}[No uniform interior gap]\label{prop:ob-gap}

There is a real-type meromorphic function \(M_*\) of order at most one such that \(M_*(1-s)=M_*(s)\), all its zeros and poles lie in \(2/5\le\Re s\le3/5\), and its right boundary profile is \(\mathbf1_{\{0\}}\), while both zeros and poles strictly inside the strip approach the line \(\Re s=3/5\) at unbounded heights.

\end{proposition}

\begin{proof} Set \(z=s-1/2\), \(a=1/10\), and, for \(n\ge1\),

\begin{equation}\label{eq:ob6-1}
u_n=a-\frac1{n+12}+in,\qquad
v_n=a-\frac2{n+24}+i(n+1/2).
\end{equation}

Their real parts lie strictly between zero and \(a\). Define

\begin{equation}\label{eq:ob6-2}
P_u(z)=\prod_{n\ge1}
\left(1-\frac{z^2}{u_n^2}\right)
\left(1-\frac{z^2}{\overline u_n^2}\right),
\end{equation}

and define \(P_v\) in the same way. These products converge uniformly on compact sets since \(\sum|u_n|^{-2}\) and \(\sum|v_n|^{-2}\) converge. Moreover

\begin{equation}\label{eq:ob6-3}
|P_u(z)|, |P_v(z)|
\le\prod_{n\ge1}(1+|z|^2/n^2)^2
=\left(\frac{\sinh(\pi|z|)}{\pi|z|}\right)^2.
\end{equation}

Thus they are entire of exponential type and order at most one. Put

\begin{equation}\label{eq:ob6-4}
M_*(s)=
\frac{P_v(s-1/2)}{
((s-1/2)^2-a^2)P_u(s-1/2)}.
\end{equation}

Its numerator and denominator are entire of order at most one, so it is meromorphic of order at most one. Both products are even and of real type, giving the required symmetries. All listed zeros and poles are distinct: the nonreal poles have integral imaginary parts, and the nonreal zeros have half-integral imaginary parts. There are simple real poles at \(s=3/5,2/5\). The right-hand interior poles have coordinates

\begin{equation}\label{eq:ob6-5}
s=\frac35-\frac1{n+12}+in,
\end{equation}

and the corresponding zeros have coordinates

\begin{equation}\label{eq:ob6-6}
s=\frac35-\frac2{n+24}+i(n+1/2).
\end{equation}

Their conjugates and reflected points complete the divisor. This proves every asserted property. For real \(\sigma>3/5\), the conjugate factor pairs are positive and \((\sigma-1/2)^2-a^2>0\), so \(M_*(\sigma)>0\) as well. \end{proof}

No Dirichlet expansion or Euler product is claimed for \(M_*\). Its role is exact: the reflection equation, finite order, local finiteness, a positive simple real pole, and the permitted boundary profile do \emph{not} imply a uniform distance between the remaining divisor and the outer line.

\subsection{A reflected positive Dirichlet example without bounded traces}

For a complementary counterexample, let \(\chi_5\) be the even primitive quadratic character modulo five and put

\begin{equation}\label{eq:ob7-1}
w=3s-1,\qquad
Q_0(s)=\frac{\zeta(w)}{L(w,\chi_5)},\qquad
M_0(s)=5^{-w/2}Q_0(s).
\end{equation}

The completed degree-one functions have the same gamma factor and root number \(+1\). Since \(w(1-s)=1-w(s)\), their completed ratio satisfies

\begin{equation}\label{eq:ob7-2}
M_0(1-s)=M_0(s).
\end{equation}

The standard zeta and primitive Dirichlet functional equations used here are recorded in \citet[\S\S25.4 and 25.15, equations 25.4.1 and 25.15.5]{dlmf}. For this particular character, its Gauss sum is

\begin{equation}\label{eq:ob7-3}
2\cos(2\pi/5)-2\cos(4\pi/5)=\sqrt5,
\end{equation}

which checks the root number directly. The ratio is globally meromorphic of order at most one.

The Euler factors of \(Q_0\) are

\begin{equation}\label{eq:ob7-4}
\begin{cases}
(1-5^{1-3s})^{-1},&p=5,\\
(1+p^{1-3s})/(1-p^{1-3s}),&\chi_5(p)=-1,\\
1,&\chi_5(p)=1.
\end{cases}
\end{equation}

They give an ordinary Dirichlet series with nonnegative integer coefficients, supported on cubes, and a logarithmic derivative with nonnegative coefficients. It converges at \(s=1\), has abscissa \(2/3\), and has a simple real pole there because \(\zeta(w)\) has its simple pole at \(w=1\) while \(L(1,\chi_5)>0\). The last inequality follows directly by grouping the convergent series into blocks:
\[
L(1,\chi_5)=\sum_{n\ge0}
\left(\frac1{(5n+1)(5n+2)}
-\frac1{(5n+3)(5n+4)}\right)>0.
\]
Thus positivity, global order one, reflection, and convergence at one alone do not rule out an off-central real pole.

The crucial arithmetic bound fails explicitly. At an inert prime \(p\), the coefficient of \(p^{-3ks}\) in \(-Q_0'/Q_0\), for odd \(k\), is

\begin{equation}\label{eq:ob7-5}
6p^k\log p.
\end{equation}

The normalized prime-power traces are unbounded. For the fixed inert prime \(p=2\), the ordinary coefficient at \(2^{3k}\) is \(2^{k+1}\). Thus the Ramanujan estimate \(a(n)\ll_\epsilon n^\epsilon\) fails for every fixed \(\epsilon<1/3\), without needing any assertion about the distribution of inert primes. This example is therefore excluded by \eqref{eq:bound}. The unscaled ratio \(\zeta(s)/L(s,\chi_5)\), already discussed in Proposition~\ref{prop:counterexample}, has its pole at one; the substitution here moves the abscissa below one while preserving reflection about \(1/2\), at exactly the cost displayed in \eqref{eq:ob7-5}.

\paragraph{Joint scope.}
No example in this appendix satisfies the complete arithmetic comparison package while violating the conclusion of Theorem~\ref{thm:main}. Proposition~\ref{prop:ob-primes} preserves the exact prime Euler factors and bounded traces but fails reflection. Proposition~\ref{prop:ob-gap} preserves reflection and order one but asserts no Euler structure. The scaled classical ratio has reflection and positivity but fails the local coefficient bound. The unresolved problem is to use the global functional equation and bounded prime Euler structure jointly to control the interior divisor. The obstruction to canonical positive removal is not an obstruction to every possible analytic reduction.
